\chapter{Main Function}
\section{Main Function}

The \textbf{Airport License Plate Detection System} is a smart camera-based solution that automatically recognises and records the license plates of vehicles entering and leaving the airport area.  
It is specially designed for \textbf{German license plates} and helps manage short-term parking at the airport gate efficiently and fairly.  
The device uses the compact and reliable \textbf{NVIDIA Jetson Nano} computer, which performs all detection, timing, and billing functions locally — without requiring any internet or cloud connection \cite{nvidia2024jetson}.

\vspace{0.5em}
When a vehicle enters, the camera captures the license plate and stores it together with the date and time.  
When the same vehicle exits, the system detects it again and automatically calculates how long it stayed inside the airport area.  
If the duration is more than \textbf{15 minutes}, the system marks it for billing.  
If the stay is \textbf{15 minutes or less}, the parking is free of charge \cite{pathak2023edgeai}.

\vspace{0.5em}
Certain vehicles such as \textbf{emergency services (ambulance, police, fire brigade)} and \textbf{official government vehicles} are automatically identified and exempted from all charges.  
These special plates are pre-registered in the system’s database, ensuring that authorised vehicles always pass freely without manual intervention.

\vspace{0.5em}
All vehicle data, including plate number, entry time, exit time, duration, and payment status, are securely stored inside the device.  
This ensures transparency, traceability, and smooth traffic flow at both the entry and exit points of the airport \cite{ultralytics2023yolo, jaided2021easyocr}.

\vspace{0.3em}
\noindent\textbf{Key Highlights:}
\begin{itemize}
	\item Recognises and records German license plates automatically.
	\item Calculates parking duration and applies the 15-minute free rule.
	\item Automatically exempts emergency and government vehicles from charges.
	\item Works fully offline — no internet or network setup required.
	\item Stores all records locally for reporting, auditing, or payment.
\end{itemize}


% --------------------------------------------------
% Section 1: How the System Works
% --------------------------------------------------
\clearpage
\section{How the System Works}

\noindent
\begin{minipage}{\linewidth}
	\centering
	\vspace{-0.2cm}
	\begin{tikzpicture}[
		scale=0.9, transform shape,
		font=\small\sffamily,
		node distance=1.0cm and 1.8cm,
		>=latex, thick,
		startstop/.style={rectangle, rounded corners=6pt, draw=gray!80, fill=gray!10,
			text centered, minimum width=3.3cm, minimum height=0.9cm},
		process/.style={rectangle, rounded corners=6pt, draw=blue!70!black, fill=blue!10,
			text centered, minimum width=3.8cm, minimum height=0.9cm},
		decision/.style={diamond, draw=orange!80!black, fill=orange!10,
			text width=3.1cm, align=center, aspect=1.7, inner sep=1pt},
		terminator/.style={ellipse, draw=gray!70!black, fill=gray!15,
			minimum width=3.3cm, text centered},
		arrow/.style={->, thick, draw=black!70},
		every node/.style={align=center}
		]
		
		% ---- Nodes (Vertical layout) ----
		\node[startstop] (start) {\textbf{System Start}};
		\node[process, below=of start] (entry) {\faCar~Vehicle Entry\\Plate Captured};
		\node[process, below=of entry] (store) {\faSave~Save Entry Time};
		\node[decision, below=of store] (exempt) {Emergency or\\Government Vehicle?};
		\node[process, right=5.2cm of exempt] (exfree) {\faCheckCircle~Free Exit\\(Exempt Vehicle)};
		\node[process, below=of exempt] (exit) {\faCarSide~Vehicle Exit\\Plate Detected Again};
		\node[process, below=of exit] (duration) {\faClock~Calculate Stay Duration};
		\node[decision, below=of duration] (check15) {Stay $\leq$ 15 minutes?};
		\node[process, right=3.4cm of check15] (free) {\faSmile~Free Exit\\(Short Stay)};
		\node[process, below=of check15] (charge) {\faBalanceScale~Display the amount user needs to Pay};
		\node[process, below=of charge] (log) {\faDatabase~Save Record\\(Billing / Audit)};
		\node[terminator, below=of log] (end) {\textbf{End of Process}};
		
		% ---- Arrows ----
		\draw[arrow] (start) -- (entry);
		\draw[arrow] (entry) -- (store);
		\draw[arrow] (store) -- (exempt);
		\draw[arrow] (exempt) -- node[above]{Yes} (exfree);
		\draw[arrow] (exempt) -- node[left]{No} (exit);
		\draw[arrow] (exit) -- (duration);
		\draw[arrow] (duration) -- (check15);
		\draw[arrow] (check15) -- node[above]{Yes} (free);
		\draw[arrow] (check15) -- node[left]{No} (charge);
		\draw[arrow] (charge) -- (log);
		\draw[arrow] (free) |- (log);
		\draw[arrow] (exfree) |- (log);
		\draw[arrow] (log) -- (end);
		
	\end{tikzpicture}
	\vspace{-0.2cm}
	\captionof{figure}{Final airport entry–exit workflow with exemptions, duration check, and billing logic.}
	\label{fig:airport_flow_vertical}
\end{minipage}


\vspace{0.3cm}
\noindent
The system repeats this workflow continuously throughout the day.  
By automatically checking for exempt vehicles and applying the 15-minute free rule, it guarantees fairness, speed, and efficiency for all users.

% --------------------------------------------------
% Section 2: Operating Environment
% --------------------------------------------------

\begin{figure}[H]
	\centering
	\fbox{
		\begin{minipage}[c][5cm][c]{0.8\textwidth}
			\centering
			\textbf{[Future Image Placeholder]}\\[4pt]
			Concept image of the airport setup will be added here.
		\end{minipage}
	}
	\caption{Placeholder: Concept image showing system deployment.}
	\label{fig:future_concept}
\end{figure}

The system is built to work smoothly and automatically with no special setup.  
Once the device is connected to power, it starts on its own and is ready for use within a few seconds.  
The camera watches both the entry and exit lanes, while the screen shows every vehicle that passes — along with how long it stayed and whether any charge applies.

\vspace{0.4em}
All processing happens inside the compact Jetson Nano computer.  
There is no need for an internet connection or external software.  
The system is designed to be reliable even during continuous, all-day operation.

\vspace{0.8em}
\noindent
\textbf{Camera Placement and Setup:}
\begin{itemize}
	\item Mount the camera about \textbf{2.0–2.5 metres above ground level}, facing the vehicle lane at a \textbf{30–40° angle}.
	\item Keep a distance of about \textbf{3–5 metres} between the camera and the point where vehicles stop or slow down.
	\item Avoid placing the camera directly opposite bright lights or reflective surfaces.
	\item Ensure there is enough ambient light. During night hours, add a small \textbf{LED lamp or floodlight} near the gate for better clarity.
	\item Keep the camera lens clean and free of dust, water spots, or scratches.
	\item Do not position the camera behind glass or near strong sunlight to prevent reflections.
\end{itemize}

\vspace{0.8em}
\noindent
\textbf{Plausibility and Manual Verification:}
\begin{itemize}
	\item If the system cannot read a plate clearly (for example, due to dirt, glare, or motion blur), it will display a \textbf{“Detection Unclear”} message on the screen.
	\item In such cases, the operator can manually enter the license plate number using a simple on-screen input field or log sheet.
	\item The system checks for \textbf{duplicate or incomplete entries} to avoid double billing.
	\item If an entry or exit is missing, the software automatically flags the record for review.
	\item For fairness, vehicles with unreadable plates are temporarily marked as \textbf{“Pending Verification”} until corrected by staff.
	\item Regular cleaning of the lens and maintaining good lighting reduces the chance of detection errors.
\end{itemize}

\vspace{0.8em}
\noindent
\textbf{In simple terms:}
\begin{itemize}
	\item Plug in the power — the system starts automatically.
	\item The camera reads license plates at entry and exit points.
	\item The display shows vehicle timing and free/charge status.
	\item Emergency or government vehicles are recognised and never charged.
	\item In rare cases of unclear detection, staff can correct the record manually.
	\item Everything runs locally — no internet or cloud connection required.
\end{itemize}

---

% --------------------------------------------------
% Section 3: User Interaction
% --------------------------------------------------
\section{User Interaction}

The system is designed for daily airport use with minimal user effort.  
All detections, timings, and charge decisions appear automatically on the display — no keyboard or technical input is required.  
The operator only monitors the live feed and confirms that vehicles are captured correctly.

\vspace{0.6em}
\noindent
\textbf{Interface Overview:}
\begin{itemize}
	\item The screen displays the \textbf{camera view} on the left and the \textbf{detection log} on the right.
	\item Each entry shows the plate number, time, duration, and parking status — “Free” (green) or “Charge” (orange).
	\item Exempt vehicles such as police or ambulance appear in blue with “No Charge.”
\end{itemize}

\vspace{0.6em}
\noindent
\textbf{User Feedback and Manual Input:}
\begin{itemize}
	\item A short sound confirms each successful detection.
	\item If a plate is unclear, the alert \textit{“Detection Unclear – Please Verify”} appears.
	\item The operator can type the correct number manually; it is saved instantly as a verified record.
\end{itemize}

\vspace{0.6em}
\noindent
\textbf{Managing Exempt Vehicles:}
\begin{itemize}
	\item Exempt vehicle numbers are stored in a simple text file (\texttt{exempt\_list.txt}).
	\item Staff can edit this list before startup; changes take effect immediately.
\end{itemize}

\vspace{0.6em}
\noindent
\textbf{Daily Operation:}
\begin{itemize}
	\item Power on — the system starts automatically.
	\item Observe the live feed and alerts.
	\item Update exempt list when needed.
	\item Turn off display and power after working hours; all data is saved automatically.
\end{itemize}

\vspace{0.4cm}
\noindent
This workflow ensures clear, automatic, and reliable operation for non-technical staff, focusing on smooth vehicle flow and accurate billing.